% ============================================
% Title Options
% ============================================

\title{Interspecies Interactions Drive Community-Level Selection in Microbial Coalescence}



% Alternative title:
% Interaction strength determines when microbial communities behave as units of selection

% ============================================
% Authors and Affiliations
% ============================================
\author[1]{Jinyeop Song}
\author[1]{Jiliang Hu}
\author[1]{Jeff Gore}

\affil[1]{Department of Physics, Massachusetts Institute of Technology, Cambridge, MA, USA}

\date{}

\maketitle

% ============================================
% Abstract
% ============================================
\begin{abstract}
It has long been debated in ecology whether communities behave as cohesive units or as loose collections of independent species. Here, we study this question in the context of community coalescence, the mixing of previously isolated communities, using bacterial microcosm experiments combined with ecological modeling. Our results demonstrate that interspecies interaction strength determines whether communities or species are the units of selection during coalescence. When interactions are moderate to strong, one parental community consistently outcompetes the other, indicating community-level selection. In contrast, under weak interactions, species fates are uncorrelated and the two communities contribute equally to the coalesced outcome, indicating the absence of community-level selection. These patterns extend to communities derived from natural samples with greater taxonomic diversity and richness. Furthermore, we identify two distinct regimes underlying community-level selection in experiments with different media conditions: an emergent regime in which collective dynamics shape outcomes that cannot be predicted from species traits alone, and a top-down regime where dominant species determine the winning community. Together, these results reconcile conflicting observations on community-level selection during community coalescence by demonstrating that communities behave as cohesive units only when interactions are sufficiently strong.
\end{abstract}

\newpage
